% Part: first-order-logic
% Chapter: tableaux
% Section: rules-and-proofs

\documentclass[../../../include/open-logic-section]{subfiles}

\begin{document}

\iftag{FOL}
      {\olfileid{fol}{tab}{rul}}
      {\olfileid{pl}{tab}{rul}}

\olsection{বিধি ও \usetoken{P}{tableau}}

!!^a{tableau} হলো কোনো !!a{structure}-এ একটি !!a{sentence} কী কী
উপায়ে সত্য বা মিথ্যা হতে পারে তার নিয়মবদ্ধ পর্যালোচনা। ট্যাবলোর
গঠন-উপাদান হলো !!{signed formula}s: একটি !!{sentence}s-এর সঙ্গে
একটি সত্যমানের “চিহ্ন”—হয় $\True$, নয়তো~$\False$। এই
!!{formula}s নিচের দিকে বাড়তে থাকা একটি বৃক্ষে সাজানো হয়।

\begin{defn}
  একটি \emph{!!{signed formula}} হলো একটি সত্যমান ও একটি
  !!a{sentence}-এর জোড়া, অর্থাৎ নিচের যেকোনো একটি:
  \[
  \sFmla{\True}{!A} \text{ অথবা } \sFmla{\False}{!A}.
  \]
\end{defn}

সহজভাবে, কোনো !!{structure}-এ $\sFmla{\True}{!A}$-কে “$!A$ সত্য
হতে পারে” এবং $\sFmla{\False}{!A}$-কে “$!A$ মিথ্যা হতে পারে” বলে
পড়া যায়।

বৃক্ষের প্রতিটি !!{signed formula} হয় একটি \emph{অনুমিতি}
(এগুলি বৃক্ষের একেবারে ওপরে লেখা থাকে), নয়তো তার ওপরের কোনো
!!a{signed formula}-এ কয়েকটি অনুমান-বিধির একটির প্রয়োগে পাওয়া যায়।
আগের !!{formula}-টির সম্ভাব্য প্রতিটি !!{main operator}-এর জন্য দুটি
বিধি আছে—চিহ্ন~$\True$ হলে একটি এবং চিহ্ন~$\False$ হলে একটি। কোনো
কোনো বিধি বৃক্ষকে শাখায় ভাগ করে, আর কোনো বিধি কেবল শাখায়
!!{signed formula}s যোগ করে। কোনো বিধি তার ঠিক আগের
!!{signed formula}-তেই প্রয়োগ করতে হবে এমন নয়; বৃক্ষের মূল থেকে বিধি
প্রয়োগের স্থান পর্যন্ত শাখার যেকোনো !!{signed formula}-এ বিধিটি
প্রয়োগ করা যায় (এবং প্রায়ই তা করতেই হয়)।

কোনো শাখায় $\sFmla{\True}{!A}$ ও $\sFmla{\False}{!A}$ উভয়ই থাকলে
শাখাটি \emph{বদ্ধ}। যে !!{tableau}-এর প্রত্যেকটি শাখা বদ্ধ, সেটি
বদ্ধ ট্যাবলো। সহজ ব্যাখ্যায় প্রতিটি শাখা একটি যৌথ সম্ভাবনা
বর্ণনা করে, কিন্তু $\sFmla{\True}{!A}$ ও $\sFmla{\False}{!A}$ একই
সঙ্গে সম্ভব নয়। অর্থাৎ কোনো শাখা বদ্ধ হলে তার বর্ণিত সম্ভাবনাটি
বাদ পড়ে যায়। বিশেষত, বদ্ধ !!{tableau} এমন সব সম্ভাবনাই বাদ দেয়
যেখানে $\sFmla{\True}{!A}$ রূপের প্রত্যেক অনুমিতিকে একই সঙ্গে সত্য
এবং $\sFmla{\False}{!A}$ রূপের প্রত্যেক অনুমিতিকে একই সঙ্গে মিথ্যা
করা যায়।

$!A$-এর \emph{জন্য বদ্ধ !!{tableau}} হলো এমন একটি বদ্ধ
!!{tableau} যার মূল~$\sFmla{\False}{!A}$। এমন বদ্ধ !!{tableau}
থাকলে $!A$-এর মিথ্যা হওয়ার সব সম্ভাবনা বাদ পড়ে; অর্থাৎ প্রতিটি
!!{structure}-এ $!A$ সত্য হতেই হবে।

\end{document}
